\subsection{FUNCTIONS}

DEFINITION 9. A graph F is said to be a functional graph if for each x there is at most one object which corresponds to x under F (Chapter I, §5, no. 3). A correspondence \(f = (\mathbf{F}, \mathbf{A}, \mathbf{B})\) is said to be a function if its graph F is a functional graph and if its source A is equal to its domain \(\mathbf{pr}_1\mathbf{F}\) . In other words, a correspondence \(f = (\mathbf{F}, \mathbf{A}, \mathbf{B})\) is a function if for every x belonging to the source A of f the relation \((x, y) \in \mathbf{F}\) is functional in y (Chapter I, §5, no. 3); the unique object which corresponds to x under f is called the value of f at the element x of A, and is denoted by \(f(x)\) (or \(f_x\) , or \(\mathbf{F}(x)\) , or \(\mathbf{F}_x\) ).

If \(f\) is a function, \(F\) its graph, and \(x\) an element of the domain of \(f\) , the relation \(y = f(x)\) is then equivalent to \((x, y) \in F\) (Chapter I, §5, no. 3, criterion C46).

Remark. The reader should beware of the confusion which may arise from the simultaneous use of the notations \(f(x)\) and \(f(X)\) (synonymous with \(f\langle X\rangle\) ) introduced in Definition 3 (cf.~Exercise 11).

Let A and B be two sets; a mapping of A into B is a function f whose source (which is equal to its domain) is equal to A and whose target is equal to B; such a function is also said to be defined on A and to take its values in B.

Instead of the phrase ``let f be a mapping of A into B'', the following phrases are often used: ``let \(f: A \rightarrow B\) be a mapping'' or even ``let \(f: A \rightarrow B\) ''. To simplify the presentation of an argument involving several mappings, we use diagrams such as

\[
\mathrm{A} \xrightarrow {f} \mathrm{B} \xrightarrow [ h ]{g} \mathrm{C} \xrightarrow [ j ]{i} \mathrm{E}
\]

in which a group of signs such as \(\mathbf{A} \to \mathbf{B}\) is to be interpreted as meaning that \(f\) is a mapping of \(\mathbf{A}\) into \(\mathbf{B}\) .

A function \(f\) defined on \(\mathbf{A}\) is said to transform \(x\) into \(f(x)\) (for all \(x \in \mathbf{A}\) ); \(f(x)\) is called the transform of \(x\) by \(f\) or (by abuse of language) the image of \(x\) under \(f\) .

Under certain circumstances, a functional graph is called a family; the domain is then called the index set, and the range is called (by abuse of language) the set of elements of the family. It is mainly in this connection that the indicial notation \(f_{x}\) is used to denote the value of f at the element x. When the index set is the product of two sets, we often speak of a double family.

Likewise a function whose target is E is sometimes called a family of elements of E. If every element of E is a subset of a set F, we speak of a family of subsets of F.

Throughout this series we shall often use the word ``function'' in place of ``functional graph''.

\minorheading{Examples of functions}

\begin{enumerate}
\def\labelenumi{(\arabic{enumi})}
\item
  The empty set is a functional graph. Every function whose graph is empty has domain and range equal to the empty set. Among such functions, the one whose target is empty (i.e., the function \((\phi, \phi, \phi)\) ) is called the empty function.
\item
  Let A be a set. The identity correspondence of A (no. 3) is a function, called the identity mapping of A.
\end{enumerate}

Thus with every set A there is associated a family, defined by the identity mapping of A, whose index set is A and whose set of elements is A. By abuse of language, a set is sometimes referred to as a ``family'', in which case it is the family thus associated with the set in question.

A function \(f\) is said to be constant if for all \(x\) and \(x'\) in the domain of \(f\) we have \(f(x) = f(x')\) .

Let \(f\) be a mapping of a set \(\mathbf{E}\) into \(\mathbf{E}\) . An element \(x\) of \(\mathbf{E}\) is said to be fixed under \(f\) if \(f(x) = x\) .

